% ==================================================================
\section{Proposed Experimental Validation}
\label{sec:future}
% ==================================================================

The central claims of this paper can be reduced to two testable
propositions. First, that LLMs inherit the positive bias of the
scientific literature. Second, that structured failure data corrects
this bias. We outline one experiment for each.

\subsection{Experiment 1. Do LLMs Overestimate Success?}

\textbf{Hypothesis.} LLMs systematically overestimate the
probability of experimental success due to positive publication bias
in their training data.

\textbf{Design.} Select 200--500 completed clinical trials from
ClinicalTrials.gov. For each trial, extract the information available
at the start of the trial (hypothesis, drug, target disease, phase,
design) and withhold the actual outcome. Prompt multiple LLMs to
estimate the probability of success based on the starting information
alone. Compare LLM predictions against actual outcomes.

\textbf{Key measures.} Overall AUC-ROC. Distribution of predicted
success probabilities for trials that actually failed. Comparison
across trial phases.

\textbf{Expected finding.} LLMs will assign systematically higher
success probabilities to trials that ultimately failed, reflecting
the positive skew of the literature they were trained on.

\subsection{Experiment 2. Does Failure Data Help?}

\textbf{Hypothesis.} Incorporating failure data into LLM context
improves predictive accuracy, but only when failure data is properly
classified.

\textbf{Design.} Using the same clinical trial dataset, compare LLM
performance under five conditions.

\begin{table}[h]
\centering
\begin{tabular}{ll}
\toprule
\textbf{Condition} & \textbf{Context provided} \\
\midrule
A (baseline)     & No additional context \\
B (success only) & Related successful trials only \\
C (failure only) & Related failed trials only \\
D (both)         & Related successful and failed trials \\
E (classified)   & D + failure cause analysis (Type A/B/C) \\
\bottomrule
\end{tabular}
\end{table}

\textbf{Key measures.} Prediction AUC-ROC per condition.
Dose-response curve (10, 50, 100, 200 failure records).

\textbf{Expected finding.} Condition B (success only) will
\textit{worsen} performance relative to baseline A, by reinforcing
positive bias. Condition D will outperform C, and E will outperform
D. If unclassified failure data degrades performance relative to
classified data, this validates the necessity of the proposed
taxonomy.

% ==================================================================
\section{Discussion}
\label{sec:discussion}
% ==================================================================

\subsection{Not All Failures Are Equal}

If failure data is to be published and used, a natural question
arises about whether all failures carry equal informational value.
We believe they do not, and that failing to distinguish between
types of failure risks degrading rather than improving any system
that consumes the data.

Consider a rejected paper. It may have been rejected because its
methodology was flawed, or because its hypothesis was genuinely
unsupported by a well-designed experiment. The former teaches us
\textit{how not to do science}. The latter teaches us \textit{what
nature does not permit}. A third category exists where methodology
and outcome are entangled, making it unclear whether the null
finding reflects a true absence of effect or a flaw in execution.

We tentatively propose a three-part classification. Type A
(methodological failure) captures cases where the approach itself
was flawed, such as applying an inappropriate statistical test,
using an underpowered sample, or failing to control for known
confounders. Type B (substantive null result) captures cases where
a sound experiment yielded a negative outcome, such as a
well-powered clinical trial that finds no effect or a carefully
controlled replication that fails to reproduce a prior finding.
Type C (ambiguous) captures cases where the two cannot be
separated, as when a reviewer notes that it is unclear whether the
null result reflects a true absence of effect or inadequate
execution.

\subsection{What Counts as a Meaningful Failure Publication}

This paper is not an argument for publishing every unsuccessful
attempt or minor negative variation. An uncurated flood of
low-effort failure reports would introduce its own form of noise,
potentially worse than the current silence. The value of failure
data depends entirely on how it is defined, structured, and
quality-controlled.

What constitutes a meaningful failure publication will differ across
disciplines. A well-powered clinical trial that finds no effect is
qualitatively different from a preliminary computational experiment
that did not converge. The former represents substantial investment
and produces definitive evidence about a hypothesis. The latter may
reflect an implementation choice rather than a scientific finding.
Each field will need to develop its own criteria for what meets the
threshold of a publishable failure, just as each field has developed
its own standards for what meets the threshold of a publishable
success. This process of definition and community acceptance is a
prerequisite for any failure publication infrastructure to function,
and cannot be imposed top-down.

For instance, a substantive null result from a well-designed
experiment (Type B) is almost always informative. A methodological
failure (Type A) is informative only when the error is common
enough that documenting it would prevent others from repeating it.
The threshold for publication will naturally differ between these
categories.

\subsection{Further Research Directions}

Beyond the two core experiments proposed in
Section~\ref{sec:future}, several additional lines of investigation
would strengthen the case presented here.

One natural extension is to test whether the relationship between
publication bias and LLM error holds across fields. If the published
positive-result ratio and the actual success ratio (obtainable from
trial registries) were measured independently for multiple medical
subfields, a positive correlation between a field's publication
bias index and LLM prediction error would provide mechanistic
evidence that publication bias is a direct cause of LLM
miscalibration, not merely a co-occurring phenomenon.

Another direction concerns LLM-assisted peer review. As discussed
in Section~\ref{sec:opportunity}, LLM reviewers trained on
positively-biased literature may lack the distributional knowledge
to identify flawed work. An experiment comparing LLM review quality
with and without failure data in the context, using publicly
available reviews from OpenReview, would directly test whether
failure exposure improves the ability to detect methodological
errors and reduces hallucinated praise. Given that 21\% of ICLR
2026 reviews were flagged as AI-generated
\citep{pangram2025iclr}, the practical relevance of this question
is immediate.

\subsection{Structural versus Defensive Responses}

Current responses to the publishing and review crisis, including
LLM detection tools, usage bans, and prompt-injection
countermeasures, are necessary but insufficient. They are defensive
measures in an arms race that the detection side is unlikely to win
permanently. A complementary approach is to change the structural
conditions that make gaming the system rational. If failure had a
legitimate place in the scientific record, the incentive to disguise
it as success would diminish, and the tools trained on that record
would develop a more accurate model of scientific reality.

We emphasize that the connection between failure publication and
fraud reduction is structural, not causal. It is plausible that
legitimizing failure as a recognized scholarly output would reduce
the pressure to fabricate. However, unless hiring and promotion
committees recognize failure contributions as scholarly merit, the
underlying pressure remains unchanged regardless of whether
publishing infrastructure exists. The causal verification of this
link requires longitudinal intervention studies that are beyond the
scope of this work.

\subsection{Limitations}

This paper presents a conceptual framework and does not include
empirical validation. While the individual observations we draw
upon are well-established, the central claim that these converge
into a unified case for failure publication infrastructure remains
an argument, not a demonstrated fact.

Several specific limitations should be noted. First, our argument
that failure data improves LLM performance rests partly on analogy
with human expertise. Senior scientists benefit from accumulated
failure experience, but it does not automatically follow that the
same information, encoded as text, will produce comparable
improvements in a statistical language model. The mechanisms of
human intuition and LLM pattern matching are fundamentally
different, and the transferability of failure information across
these systems requires direct experimental verification.

Second, we do not address the practical challenges of incentivizing
failure publication at scale, including questions of intellectual
property, competitive risk, and the additional labor burden on
researchers who are already under pressure. A failure publication
system that no one uses solves nothing.

Third, the proposed failure taxonomy is preliminary and untested.
Its value depends on whether the categories can be applied
consistently and whether the distinction actually affects downstream
LLM performance, both of which are empirical questions.